% Generated from locale/tr/content/proof-theory/normalization/translations.tex; do not hand-edit.


\olfileid[tr]{pt}{nor}{tr}

\olsection{Normal \Log{N2i}-\usetoken{P}{proof} ile \Log{G2i} Arasında Çeviri}

\Log{N2i} içindeki !!^{proof} örnekleri, \Log{G2i} içindeki !!{proof} biçimine
ve bunun tersi yönde çevrilebilir. Öncelikle, kesmesiz \Log{G2i}
!!{proof} örnekleri çevrildiğinde normal \Log{N2i}-!!{proof} örnekleri elde edilir.

Sonucu ifade etmek için \Log{N2i} ve \Log{G2i} içindeki sekantların
önbileşenlerini daha kolay ilişkilendirmemizi sağlayacak biraz
terminolojiye gereksinimimiz var: Etiketli !!{formula} biçimlerinden oluşan bir
$\Gamma'$ kümesi, etiketsiz !!{formula} biçimlerinden oluşan $\Gamma$ çoklu
kümesine şu durumda \emph{karşılık gelir}: Her $!A$ !!{formula} için,
$\Gamma'$ içinde $x : !A$ biçimindeki !!{element} sayısı, $!A$'nın
$\Gamma$ içindeki katlılığından küçük ya da ona eşittir.

\begin{cor}
$\Log{G2i} \Proves \Gamma \Sequent !C$ ise $\Gamma' \Sequent !C$
için normal bir \Log{N2i}-!!{proof}~$\delta$ vardır; burada etiketli
!!{formula} biçimlerinden oluşan $\Gamma'$ kümesi $\Gamma$ çoklu kümesine
karşılık gelir. Başka bir deyişle, her $!A$ !!{formula} için,
$\Gamma'$ içinde $x : !A$ biçimindeki !!{element} sayısı $!A$'nın
$\Gamma$ içindeki katlılığından (yani $!A$ kopyalarının sayısından)
küçük ya da ona eşittir.
\end{cor}

\begin{proof}
\olref[nat][tgi]{prop:G2i-to-N2i} ispatını inceleyelim:
!!{introduction} kurallarını !!{proof} sonlarına, !!{elimination}
kurallarını ise yalnızca !!{proof} başlarına ekleriz; dolayısıyla hiçbir
!!{introduction} kuralını !!a{elimination} kuralının üzerine getirmeyiz.
\end{proof}

Ancak \olref[nat][tgi]{prop:G2i-to-N2i}'de verilen \Cut{} kuralı
çevirisi, \Log{N2i}-!!{proof}~$\delta$ içinde kesmeler ortaya
çıkarabilir. Bir \Log{G2i} !!{proof}, sırasıyla $\Gamma_1 \Sequent !A$
ve $!A, \Gamma_2 \Sequent \Delta_2$ için doğrudan alt-!!{proof} örnekleri
$\pi_1$ ve $\pi_2$ ile \Cut{} çıkarımında bitiyorsa $\pi_1$'in
$\delta_1$ çevirisini, $\pi_2$'nin $\delta_2$ çevirisinde geçen $x:!A
\Sequent !A$ aksiyomlarına aşılarız. $\delta_1$, !!a{introduction}
kuralında ana !!{formula} olarak $!A$'yı taşıyarak bitiyorsa ve
$\delta_2$ içindeki $x:!A \Sequent !A$ aksiyomu !!a{elimination}
kuralının büyük öncülüyse bunun sonucunda $\delta$ içinde bir kesme
oluşur.

\Log{N2i} içindeki !!{proof} örnekleri, \Log{G2i} içindeki !!{proof} biçimine
de çevrilebilir. \olref[nat][tni]{prop:N2-to-G2} ispatında
\Elim{\land}, \Elim{\lif}, \Elim{\lnot}, \Elim{\lforall} çevirisi,
elde edilen \Log{G2i} !!{proof} içinde \Cut{} çıkarımlarına yol açar.
Ancak başladığımız \Log{N2i}-!!{proof} normal ise bundan kaçınılabilir.

Bir \Log{N2i}-!!{proof}~$\delta$ içindeki sekant oluşumlarından
oluşan $S_1$, \dots,~$S_n$ listesine, $S_1$ bir aksiyomsa ve $S_i$ bir
çıkarımın büyük öncülü olduğunda $S_{i+1}$ o çıkarımın sonucuysa
\emph{dal} diyelim. Başka bir deyişle dal, $\delta$ içinde sekantları
aksiyomlardan aşağı doğru izler, ancak !!{elimination}
kurallarının küçük öncüllerinde durur. $\delta$ için bir \emph{ana dal},
$S_n$'nin son sekant olduğu bir daldır. Her kuralın en az bir büyük
öncülü bulunduğundan (yalnızca \Intro{\land} iki büyük öncüle sahiptir),
her !!{proof}~$\delta$ en az bir ana dala sahip olmalıdır.


\begin{prop}\ollabel{prop:normal-N2i-to-G2i} $\delta$, $\Gamma \Sequent
!A$ için normal bir $\Log{N2c}$- ya da $\Log{N2i}$-!!{proof} ise
$\Gamma' \Sequent \Delta$ için (kesmesiz) bir $\Log{G2c}$-!!{proof}
($\Log{G2i}$-!!{proof})~$\pi$ vardır; burada $\Gamma'$, $\Gamma$
içindeki etiketlerin kaldırılmasıyla elde edilen !!{formula} çoklu
kümesidir ve $!A$, $\lfalse$ değilse $\Delta = \{!A\}$, öyleyse
$\Delta = \emptyset$ olur.
\end{prop}

\begin{proof}
Bu kez tümevarım, $\Gamma \Sequent !A$ için !!{proof}~$\delta$ içindeki
çıkarımların sayısı üzerindendir. $\delta$ içinde çıkarım yoksa
$x:!A \Sequent !A$ bir aksiyomdur ve $\pi$, karşılık gelen $!A \Sequent
!A$ aksiyomudur; $!A \ident \lfalse$ ise $\lfalse \Sequent \ $ olur.

$\delta$ bir çıkarım kuralında bitiyorsa durumları ayırırız:
\begin{enumerate}
  \item $R$, !!a{introduction} kuralı ya da \FalseInt ise çeviri
  tam olarak \olref[nat][tni]{prop:N2-to-G2} ispatındaki gibi ilerler
  ve \Cut{} gerektirmez.

  $\delta$'nın son kuralı $R$, !!a{elimination} kuralıysa şunlara
  dikkat edin:
  \begin{enumerate}
    \item $\delta$'nın bir ana dalı hiçbir !!{introduction} kuralından
    geçmez; aksi hâlde ana dal boyunca !!a{introduction} kuralını ya da
    \FalseInt{} kuralını !!a{elimination} kuralı izler ve $\delta$
    normal olmazdı.
    \item $\delta$'nın ana dalları boyunca !!{introduction} kuralı
    bulunmadığından yalnızca bir ana dal vardır. (Yalnızca
    \Intro{\land} iki büyük öncüle sahiptir; dolayısıyla birden çok ana
    dalın bir \Intro{\land} kuralının öncüllerinden geçmesi gerekirdi.)
    \item Ana dal \Elim{\lor} ya da \Elim{\lexists} çıkarımının büyük
    öncülünü içeriyorsa bu tür yalnızca bir kural vardır ve bu da son
    kural, yani $R$'dir. (Aksi hâlde ana dal, sonucu !!a{elimination}
    kuralının büyük öncülü olan bir \Elim{\lor} ya da \Elim{\lexists}
    kuralı içerirdi ve permütasyon dönüştürmesi uygulayabildiğimiz için
    ispat normal olmazdı.)
  \item \Elim{\lor} ve \Elim{\lexists} dışındaki hiçbir
  !!{elimination} kuralı varsayımları boşa çıkarmaz, yani sekantların
  sol bağlamını değiştirmez. \Elim{\lor} ve \Elim{\lexists}
  yalnızca küçük bir öncülün üzerindeki varsayımları boşa çıkarır; yani
  yalnızca küçük öncüllerinin sol bağlamını değiştirir. Başka bir
  deyişle, $\delta$'nın ana dalındaki bir $S_i$ sekantı
  $\Gamma_1$ bağlamına sahipse $S_i$'nin büyük öncül olduğu çıkarımın
  $S_{i+1}$ sonucu $\Gamma_1' \supseteq \Gamma_1$ bağlamına sahiptir.
  \item Dolayısıyla $x: !C \Sequent !C$, ana daldaki en üst sekant
  $S_1$ ise $x:!C \in \Gamma$ olur.
  \end{enumerate}
  Şimdi $!C$'nin biçimine göre durumları ayırırız. Ana daldaki her
  $S_i$, !!a{elimination} kuralının büyük öncülü olduğundan bu,
  $x:!C \Sequent !C$ için de geçerlidir.

  \item $!C \ident !D \land !E$ olsun. O zaman $\delta$ şu biçimdedir:
  \[
  \Axiom$x:!D \land !E \fCenter !D \land !E$
  \RightLabel{\Elim{\land}[1]}
  \UnaryInf$x:!D \land !E \fCenter !D$
  \Deduce$x:!D \land !E, \Gamma_1 \fCenter !A$
  \DisplayProof
  \]
  $x:!D \land !E \Sequent !D$'yi içeren ana dal, \Intro{\lif} ya da
  \Intro{\lnot} çıkarımının hiçbir büyük öncülünü ve \Elim{\lor} ya da
  \Elim{\lexists} çıkarımının hiçbir küçük öncülünü içermediğinden,
  $x:!D \land !E$ içindeki bağlamın !!{formula} biçimleri ana dal boyunca
  değişmeden kalır. Bu, son sekantın bağlamının neden $x:!D \land
  !E$'yi içermesi gerektiğini açıklar. Ayrıca \Elim{\land} ile biten
  alt-!!{proof} yerine $x:!D \Sequent !D$ koyup ana dal boyunca her
  sekantın bağlamındaki $x:!D \land !E$ yerine $x:!D$ koyarak
  !!{proof}~$\delta$'yı bir !!{proof}~$\delta_1$ biçimine getirebiliriz;
  yani aşağıdakini
  \[
  \bottomAlignProof
  \Axiom$x:!D \land !E \fCenter !D \land !E$
  \RightLabel{\Elim{\land}[1]}
  \UnaryInf$x:!D \land !E \fCenter !D$
  \Deduce$x:!D \land !E, \Gamma_1 \fCenter !A$
  \DisplayProof
  \quad\text{şuna dönüştürürüz:}\quad
  \bottomAlignProof
  \Axiom$x:!D \fCenter !D$
  \Deduce$x:!D, \Gamma_1 \fCenter !A$
  \DisplayProof
  \]
  $\delta$ içindeki çıkarım sayısı $\delta_1$ içindeki çıkarım
  sayısından $1$ fazla olduğundan tümevarım varsayımı $\delta_1$ için
  geçerlidir.

  Tümevarım varsayımına göre $x:!D, \Gamma_1 \Sequent !A$ için bir
  \Log{G2i}-!!{proof} $\pi_1$ vardır. \LeftR{\land} uygulayarak
  !!{proof}~$\pi$'yi elde ederiz:
  \[
  \AxiomC{}
  \RightLabel{$\pi_1$}
  \Deduce$!D, \Gamma_1' \fCenter !A$
  \RightLabel{\LeftR{\land}}
  \UnaryInf$!D \land !E, \Gamma_1' \fCenter !A$
  \DisplayProof
  \]
  $!A \ident \lfalse$ ise tümevarım varsayımının verdiği $\pi_1$'in
  artbileşeni boştur; \LeftR{\land} uygulaması da artbileşeni boş bırakır.
  \item $!C \ident !D \lor !E$ ise \Elim{\lor} çıkarımının büyük
  öncülü olmalıdır ve bu çıkarım ana daldaki son çıkarım $R$ olmalıdır.
  Başka bir deyişle $\delta$ şu biçimdedir:
  \[
  \Axiom$x:!D \lor !E \fCenter !D \lor !E$
  \AxiomC{}
  \RightLabel{$\delta_1$}
  \Deduce$y:!D, \Gamma_1 \fCenter!A$
  \AxiomC{}
  \RightLabel{$\delta_2$}
  \Deduce$z:!E, \Gamma_2 \fCenter !A$
  \RightLabel{\Elim{\lor}}
  \TrinaryInf$x:!D \lor !E, \Gamma_1, \Gamma_2 \fCenter !A$
  \DisplayProof
  \]
  Tümevarım varsayımı !!{proof} örnekleri $\delta_1$ ve $\delta_2$ için
  geçerlidir; $!D, \Gamma_1' \Sequent !A$ ve $!E, \Gamma_2' \Sequent
  !A$ için \Log{G2i}-!!{proof} örnekleri $\pi_1$ ve $\pi_2$ elde ederiz.
  $\pi$'yi elde etmek için $\LeftR{\lor}$ uygularız:
  \[
  \AxiomC{}
  \RightLabel{$\pi_1$}
  \Deduce$!D, \Gamma_1' \fCenter !A$
  \AxiomC{}
  \RightLabel{$\pi_2$}
  \Deduce$!E, \Gamma_2' \fCenter !A$
  \RightLabel{\LeftR{\lor}}
  \BinaryInf$!D \lor !E, \Gamma_1', \Gamma_2' \fCenter !A$
  \DisplayProof
  \]
  $\Elim{\lor}$, $y: !D$ ya da $z: !E$'yi boşa çıkarmıyorsa (yani
  bunlar küçük öncüllerin bağlamında yoksa) birkaç
  \LeftR{\Weakening} eklememiz gerekir. $!A \ident \lfalse$ ise
  tümevarım varsayımının verdiği artbileşenler ve sonuç artbileşeni boş
  kalır. Örneğin:
  \[
  \AxiomC{}
  \RightLabel{$\pi_1$}
  \Deduce$\Gamma_1' \fCenter $
  \RightLabel{\LeftR{\Weakening}}
  \UnaryInf$!D, \Gamma_1' \fCenter $
  \AxiomC{}
  \RightLabel{$\pi_2$}
  \Deduce$\Gamma_2' \fCenter $
  \RightLabel{\LeftR{\Weakening}}
  \UnaryInf$!E, \Gamma_2' \fCenter $
  \RightLabel{\LeftR{\lor}}
  \BinaryInf$!D \lor !E, \Gamma_1', \Gamma_2' \fCenter $
  \DisplayProof
  \]
  
  \item $!C \ident !D \lif !E$ olsun. O zaman $\delta$ şu biçimdedir:
  \[
  \Axiom$x:!D \lif !E \fCenter !D \lif !E$
  \AxiomC{}
  \RightLabel{$\delta_1$}
  \Deduce$\Gamma_1 \fCenter!D$
  \RightLabel{\Elim{\lif}}
  \BinaryInf$x:!D \lif !E, \Gamma_1 \fCenter !E$
  \RightLabel{$\delta_2$}
  \Deduce$x:!D \lif !E, \Gamma_1, \Gamma_2 \fCenter !A$
  \DisplayProof
  \]
  Burada $x:!D \lif !E, \Gamma_1 \Sequent !E$'yi içeren ana dal,
  \Intro{\lif} ya da \Intro{\lnot} çıkarımının hiçbir büyük öncülünü
  ve \Elim{\lor} ya da \Elim{\lexists} çıkarımının hiçbir küçük
  öncülünü içermediğinden, $x:!D \lif !E, \Gamma_1$ içindeki bağlam
  !!{formula} biçimleri ana dal boyunca değişmeden kalır. Bu, son sekantın
  bağlamının neden $x:!D \lif !E, \Gamma_1$'i içermesi gerektiğini
  açıklar. Ayrıca \Elim{\lif} ile biten alt-!!{proof} yerine $x:!E
  \Sequent !E$ koyup ana dal boyunca her sekantın bağlamındaki
  $x:!D \lif !E, \Gamma_1$ yerine $x:!E$ koyarak !!{proof} parçası
  $\delta_2$'yi doğru bir !!{proof}~$\delta_2'$ biçimine getirebiliriz;
  yani aşağıdakini
  \[
  \bottomAlignProof
  \Axiom$x:!D \lif !E, \Gamma_1 \fCenter !E$
  \RightLabel{$\delta_2$}
  \Deduce$x:!D \lif !E, \Gamma_1, \Gamma_2 \fCenter !A$
  \DisplayProof
  \quad\text{şuna dönüştürürüz:}\quad
  \bottomAlignProof
  \Axiom$x:!E \fCenter !E$
  \RightLabel{$\delta_2'$}
  \Deduce$x:!E, \Gamma_2 \fCenter !A$
  \DisplayProof
  \]
  $\delta$ içindeki çıkarım sayısı, $\delta_1$ ve $\delta_2$ içindeki
  çıkarım sayılarının toplamına ana dalın başlangıcındaki \Elim{\lif}
  çıkarımı için $1$ eklenmesiyle elde edildiğinden, tümevarım varsayımı
  $\delta_1$ ve $\delta_2'$ için geçerlidir. (Bu çıkarım aynı zamanda
  $\delta$ içindeki son çıkarım $R$ ise $\delta_1$ sıfır çıkarım içerir.)

  Tümevarım varsayımına göre $\Gamma_1 \Sequent !D$ ve $!E, \Gamma_2'
  \Sequent !A$ için \Log{G2i}-!!{proof} örnekleri $\pi_1$ ve $\pi_2$ vardır.
  \LeftR{\lif} uygulayarak !!{proof}~$\pi$'yi elde ederiz:
  \[
  \AxiomC{}
  \RightLabel{$\pi_1$}
  \Deduce$\Gamma_1' \fCenter !D$
  \AxiomC{}
  \RightLabel{$\pi_2$}
  \Deduce$!E, \Gamma_2' \fCenter !A$
  \RightLabel{\LeftR{\lif}}
  \BinaryInf$!D \lif !E, \Gamma_1', \Gamma_2' \fCenter !A$
  \DisplayProof
  \]
  $!D \ident \lfalse$ ise $\pi_1$'i \LeftR{\lif} kuralının öncülüne
  uygun hâle getirmek için bir \RightR{\Weakening} ekleriz. Buna karşılık
  $!A \ident \lfalse$ ise $\pi_2$'nin ve elde edilen sonucun artbileşeni
  boş kalır. Örneğin her iki durum birlikte gerçekleştiğinde:
  \[
  \AxiomC{}
  \RightLabel{$\pi_1$}
  \Deduce$\Gamma_1' \fCenter $
  \RightLabel{\RightR{\Weakening}}
  \UnaryInf$\Gamma_1' \fCenter !D$
  \AxiomC{}
  \RightLabel{$\pi_2$}
  \Deduce$!E, \Gamma_2' \fCenter $
  \RightLabel{\LeftR{\lif}}
  \BinaryInf$!D \lif !E, \Gamma_1', \Gamma_2' \fCenter $
  \DisplayProof
  \]

\end{enumerate}
Kalan durumlar benzer biçimde ele alınır ve alıştırma olarak bırakılır.
\end{proof}

\begin{cor}
Her !!{formula}, normal bir \Log{N1i}- ya da \Log{N2i}-!!{proof}
içinde geçiyorsa son-!!{formula} ya da son sekantın bir alt
!!{formula} biçimidir.
\end{cor}

\begin{proof}
  \olref{prop:normal-N2i-to-G2i} uyarınca her normal
  \Log{N2i}-!!{proof}, bir \Log{G2i}-!!{proof} biçimine çevrilebilir.
  İspat incelendiğinde her !!{formula} biçiminin,
  \Log{N2i}-!!{proof} içinde geçiyorsa \Log{G2i}-!!{proof} içinde de
  geçtiği görülür.
  \Cut{} kuralını içermeyen \Log{G2i}, alt-!!{formula} özelliğine
  sahiptir.
\end{proof}

